\documentclass[11pt]{article}
\usepackage[margin=1in]{geometry}
\usepackage{times}
\usepackage{url}
\usepackage{hyperref}
\usepackage{natbib}

\title{\textbf{PersonalQuery: Reproducibility Package for Grounded Personalized Query Generation}}
\author{Anonymous\\
Anonymous Institution\\
\texttt{anonymous@email.org}}

\begin{document}
\maketitle

\begin{abstract}
This reproducibility companion presents PersonalQuery, a complete artifact package for grounded personalized query generation. The original work introduced a 12-stage methodology achieving +1.81 points improvement (45\% relative gain). Our reproducibility package includes: (1) complete 12-stage pipeline; (2) benchmark dataset with 134 paired queries across 10 users; (3) comprehensive documentation; (4) single-command reproduction scripts; and (5) artifacts publicly archived on Zenodo with DOI. The package enables full reproduction within 15-30 minutes per user with expected outputs matching paper claims within $\pm$0.2 tolerance. We target ACM badges: Artifacts Available, Artifacts Evaluated-Functional, and Results Reproduced.
\end{abstract}

\section{Introduction}

Personalized search query generation represents an emerging challenge in recommender systems research. The original PersonalQuery work~\citep{personalquery2026} introduced a novel grounded methodology extracting verified user preferences from reviews and generating queries reflecting individual writing styles.

This companion focuses on reproducibility, providing a complete artifact package enabling researchers to reproduce all main experimental results with single commands.

\section{Artifact Overview}

\textbf{Claim 1: Personalized queries significantly outperform public queries.}\\
Reproduction: Script \texttt{reproduce\_main\_results.py} evaluates query pairs and outputs scores matching Table 2 (expected: Public avg $\approx$4.0, Personalized avg $\approx$5.8, improvement $\approx$+1.81).

\textbf{Claim 2: Grounded validation filters 71.5\% of preferences.}\\
Reproduction: Pipeline execution shows preference counts before/after validation.

\textbf{Claim 3: 100\% of users show positive improvement.}\\
Reproduction: Full evaluation script outputs per-user improvements, all expected positive.

\section{Reproduction Methodology}

\textbf{Installation} (5-10 minutes):
\begin{verbatim}
git clone https://github.com/[user]/personalquery.git
cd personalquery
conda env create -f environment.yml
conda activate personalquery
python scripts/verify_setup.py
\end{verbatim}

\textbf{Main Results} (15-30 minutes):
\begin{verbatim}
python scripts/reproduce_main_results.py --user-id A13OFOB1394G31
\end{verbatim}
Expected: Public 4.0, Personalized 5.8, Improvement +1.81 (±0.2 tolerance).

\textbf{Full Evaluation} (2-4 hours):
\begin{verbatim}
python scripts/run_full_evaluation.py --all-users
\end{verbatim}

\section{Technical Requirements}

\textbf{Hardware}: CPU 8+ cores, RAM 16GB minimum, Storage 5GB.

\textbf{Software}: Python 3.8+, conda/pip, LLM API (GLM-4.5-Air or OpenAI GPT-4).

\textbf{Dependencies}: openai, spacy, sentence-transformers, pandas, numpy, scikit-learn.

\section{Artifact Availability}

\textbf{Primary}: Zenodo (\url{https://zenodo.org/record/XXXXXX})\\
\textbf{Secondary}: GitHub (\url{https://github.com/[username]/personalquery})\\
\textbf{License}: MIT (allows comparison and extension)

\bibliographystyle{plainnat}
\bibliography{references}

\end{document}
